%% THERMODYNAMIQUE %%
%
\section{Rapport de m�lange en vapeur d'eau � partir de la temp�rature du
point de ros�e}

\itemstyle{ Op�rateur~: } \operfmt{RapM\_Td}
\index{RapM\_Td}

%%%%%OBJET%%%%%
\itemstyle{ Objet~: }
Calcul du rapport de m�lange en vapeur d'eau � partir
de la temp�rature du point de ros�e.

%%%%%ARGUMENTS%%%%%
\bigskip
%%%%%ENTREES%%%%
\begin{tabarg}
\multicolumn{3}{l}{\titretab{Entr\'ees}}\\
\hline
$T_d$ & Vecteur, {\tt float} & Temp\'erature du point de ros\'ee ($K$ ou $^\circ C$)\\
$P_s$ & Vecteur, {\tt float} & Pression statique ($hPa$) \\
\end{tabarg}

\medskip
%%%%%SORTIES%%%%
\begin{tabarg}
\multicolumn{3}{l}{\titretab{Sorties}}\\
\hline
$RapM$ & Vecteur, {\tt float} & Rapport de m\'elange en vapeur d'eau ($g/kg$)\\
\end{tabarg}


%%%%%FORMULE%%%%
\itemstyle{ Formule }

%%% THERMODYNAMIQUE
\[
e = e_w(T_d) \qquad \forall T_d
\]
\[
RapM = \frac{R_a}{R_v} . \frac{e}{P_s - e } 
\]
avec $\frac{R_a}{R_v} = \frac{M_v}{M_a} = 0.622$


%%%%COMMENTAIRE%%%%
\itemstyle{Commentaire~:}

La fonction $e_{w}(t)$ est d�crite en annexe
\ref{annexe-thermo}.

Les formules employ�es correspondent � l'hypoth�se que 
l'argument d'entr�e $T_d$ est toujours un point de ros�e (et non
un point de givre si $T_d < 0$�C); 
si $T_d$ provient d'un calcul, cette hypoth�se peut �tre v�rifi�e,
mais si il provient d'un hygrom�tre � miroir
refroidi, il est probable que les valeurs en-dessous de 0�C soient
une mesure de point de givre, et il vaut mieux utiliser l'op�rateur
\hyperref[OP-RapM_Tdf]{\operfmt{RapM\_Tdf}}.


%%%%%AUTEUR%%%%%
\itemstyle{ Auteur~:}
CNRM/GMEI/TRAMM

%%%%%REFERENCES%%%%%
\itemstyle{ R�f�rences~:}
Triplet-Roche\cite{Triplet-Roche},\\
D�finitions des humidit�s (relatives, absolues, sp�cifiques)
et du rapport de m�lange
selon la table 4.3 des tables m�t�orologiques internationales
de l'OMM \cite{OMM-188.TP.94}.


